\section*{Capítulo \ref{ch:g2:measuring-length} --- Medir longitudes}

\begin{solution}{exo:g2:measuring-length:1}
Si está cerca: espalda contra espalda. Si está lejos: mídete en
\omterm{def:g2:measuring-length:units}{centímetros} y manda el número --- un número puede viajar en una
carta; una comparación espalda contra espalda, no.
\end{solution}

\begin{solution}{exo:g2:measuring-length:2}
Ninguna de las dos ha contado mal. Los pasos de Mara son más cortos,
así que caben más en la misma clase. \omterm{def:g2:measuring-length:measure}{Unidades} distintas dan números
distintos para la misma longitud.
\end{solution}

\begin{solution}{exo:g2:measuring-length:3}
El \omterm{def:g2:measuring-length:units}{centímetro} es el mismo para todo el mundo --- no es el cuerpo de
nadie. Una medida en \omterm{def:g2:measuring-length:units}{centímetros} la puede comprobar cualquier otra
persona con cualquier regla.
\end{solution}

\begin{solution}{exo:g2:measuring-length:4}
La hormiga: \unit{cm} (¡y aún menos!); la puerta: \unit{m}; el patio:
\unit{m}; el ancho de este libro: \unit{cm}.
\end{solution}

\begin{solution}{exo:g2:measuring-length:5}
No --- entre el borde de una regla y su marca del cero suele haber un
pequeño espacio en blanco, así que el lápiz es un poco más corto que
$13$ \omterm{def:g2:measuring-length:units}{centímetros}. A Tim se le olvidó poner la \emph{marca del cero},
y no el borde, en la punta del lápiz.
\end{solution}

\begin{solution}{exo:g2:measuring-length:6}
Las respuestas varían; lo habitual: la mano, unos \qty{7}{cm} de
ancho; la cuchara, unos \qty{15}{cm}; la taza, unos \qty{9}{cm}. Lo
importante: cada número lleva su \unit{cm}.
\end{solution}

\begin{solution}{exo:g2:measuring-length:7}
$100 + 100 = 200$: la puerta mide \qty{200}{cm} de alto.
\end{solution}

\begin{solution}{exo:g2:measuring-length:8}
Las respuestas varían; casi todas las camas andan cerca de
\qty{2}{m}. Lo que importa es comparar lo que se había calculado a ojo
con el número medido.
\end{solution}

\begin{solution}{exo:g2:measuring-length:9}
La cinta de \qty{45}{cm} es más larga: $45 - 38 = 7$, o sea que por
\qty{7}{cm}. Sin números, ponerlas una al lado de la otra también
descubre cuál es la más larga --- pero ``\qty{45}{cm}'' cabe en una
carta y ``una al lado de la otra'' no.
\end{solution}

\begin{solution}{exo:g2:measuring-length:10}
Nadie más tiene las piernas del abuelo: un niño pequeño que recorriera
el mismo huerto podría contar $60$ zancaditas, y una visita no podría
comprobar ``$30$ zancadas'' sin el abuelo delante. Medida en
\omterm{def:g2:measuring-length:units}{metros} --- los mismos para todo el mundo ---, la longitud del huerto
se puede escribir, mandar y comprobar por cualquiera.
\end{solution}
